\documentclass[10pt,a4paper]{amsart}
\usepackage{amsmath,amssymb,amsthm,graphicx}
% For commutative diagrams we use Xy-pic
\usepackage[all]{xy}

% We use 2cell for 2-commutative diagrams.
\xyoption{2cell}
% Theorem environments.
%
\theoremstyle{plain}
\newtheorem{theorem}[subsection]{Theorem}
\newtheorem{proposition}[subsection]{Proposition}
\newtheorem{lemma}[subsection]{Lemma}

\theoremstyle{definition}
\newtheorem{definition}[subsection]{Definition}
\newtheorem{example}[subsection]{Example}
\newtheorem{exercise}[subsection]{Exercise}
\newtheorem{situation}[subsection]{Situation}

\theoremstyle{remark}
\newtheorem{remark}[subsection]{Remark}
\newtheorem{remarks}[subsection]{Remarks}

\numberwithin{equation}{subsection}

% Macros
%
\def\lim{\mathop{\mathrm{lim}}\nolimits}
\def\colim{\mathop{\mathrm{colim}}\nolimits}
\def\Spec{\mathop{\mathrm{Spec}}}
\def\Hom{\mathop{\mathrm{Hom}}\nolimits}
\def\Ext{\mathop{\mathrm{Ext}}\nolimits}
\def\SheafHom{\mathop{\mathcal{H}\!\mathit{om}}\nolimits}
\def\SheafExt{\mathop{\mathcal{E}\!\mathit{xt}}\nolimits}
\def\Sch{\mathit{Sch}}
\def\Mor{\mathop{\mathrm{Mor}}\nolimits}
\def\Ob{\mathop{\mathrm{Ob}}\nolimits}
\def\Sh{\mathop{\mathit{Sh}}\nolimits}
\def\NL{\mathop{N\!L}\nolimits}
\def\CH{\mathop{\mathrm{CH}}\nolimits}
\def\proetale{{pro\text{-}\acute{e}tale}}
\def\etale{{\acute{e}tale}}
\def\QCoh{\mathit{QCoh}}
\def\Ker{\mathop{\mathrm{Ker}}}
\def\Im{\mathop{\mathrm{Im}}}
\def\Coker{\mathop{\mathrm{Coker}}}
\def\Coim{\mathop{\mathrm{Coim}}}

% Boxtimes
%
\DeclareMathSymbol{\boxtimes}{\mathbin}{AMSa}{"02}

%
% Macros for moduli stacks/spaces
%
\def\QCohstack{\mathcal{QC}\!\mathit{oh}}
\def\Cohstack{\mathcal{C}\!\mathit{oh}}
\def\Spacesstack{\mathcal{S}\!\mathit{paces}}
\def\Quotfunctor{\mathrm{Quot}}
\def\Hilbfunctor{\mathrm{Hilb}}
\def\Curvesstack{\mathcal{C}\!\mathit{urves}}
\def\Polarizedstack{\mathcal{P}\!\mathit{olarized}}
\def\Complexesstack{\mathcal{C}\!\mathit{omplexes}}
% \Pic is the operator that assigns to X its picard group, usage \Pic(X)
% \Picardstack_{X/B} denotes the Picard stack of X over B
% \Picardfunctor_{X/B} denotes the Picard functor of X over B
\def\Pic{\mathop{\mathrm{Pic}}\nolimits}
\def\Picardstack{\mathcal{P}\!\mathit{ic}}
\def\Picardfunctor{\mathrm{Pic}}
\def\Deformationcategory{\mathcal{D}\!\mathit{ef}}
\usepackage{cleveref}
\crefname{originallemma}{original lemma}{original lemmas}
\Crefname{originallemma}{Original lemma}{Original lemmas}
\makeatletter
\newlabel{lemma-additivity-chern-classes}{{02UI}{source}{}{sourcepointer-1}{}}
\newlabel{lemma-additivity-chern-classes@cref}{{[originallemma][0][]02UI}{[1][0][]source}{}{}{}}
\makeatother
\crefname{originallemma}{original lemma}{original lemmas}
\Crefname{originallemma}{Original lemma}{Original lemmas}
\makeatletter
\newlabel{lemma-chern-character-additive}{{0F9C}{source}{}{sourcepointer-2}{}}
\newlabel{lemma-chern-character-additive@cref}{{[originallemma][0][]0F9C}{[1][0][]source}{}{}{}}
\makeatother
\crefname{originallemma}{original lemma}{original lemmas}
\Crefname{originallemma}{Original lemma}{Original lemmas}
\makeatletter
\newlabel{lemma-chern-classes-defined}{{0GUE}{source}{}{sourcepointer-3}{}}
\newlabel{lemma-chern-classes-defined@cref}{{[originallemma][0][]0GUE}{[1][0][]source}{}{}{}}
\makeatother
\crefname{originallemma}{original lemma}{original lemmas}
\Crefname{originallemma}{Original lemma}{Original lemmas}
\makeatletter
\newlabel{perfect-proposition-perfect-resolution-property}{{0F8I}{source}{}{sourcepointer-4}{}}
\newlabel{perfect-proposition-perfect-resolution-property@cref}{{[originallemma][0][]0F8I}{[1][0][]source}{}{}{}}
\makeatother
\crefname{originallemma}{original lemma}{original lemmas}
\Crefname{originallemma}{Original lemma}{Original lemmas}
\makeatletter
\newlabel{situation-setup}{{02QL}{source}{}{sourcepointer-5}{}}
\newlabel{situation-setup@cref}{{[originallemma][0][]02QL}{[1][0][]source}{}{}{}}
\makeatother
\title{Additivity of the Chern class, Chern character and Chern polynomials along a distinguished triangle of perfect objects}
\author{The Stacks Project Authors}
\date{stacks-project:chow.tex:local-snapshot-20261012; GNU FDL 1.2 or later}
\begin{document}
\maketitle
\noindent\textit{Editorial note.} All delivered mathematical text is the author's own, copied from the pinned Stacks Project chapter file chow.tex; no statement or proof step was rewritten, reordered or removed.\par\medskip
\noindent\textit{Editorial note.} Preamble provenance: the pinned chapter file opens with \textbackslash{}input\{preamble\} and carries no preamble text of its own, so the project's single Stacks-wide preamble staged at sources/stacks-chow/preamble.tex is excerpted verbatim as two spans, lines 14-18 (Xy-pic setup) and lines 154-217 (the native theorem declarations and mathematical macros). Nothing was added, removed or redefined; the wrapper class and the amsmath/amssymb/amsthm/graphicx base packages are supplied by the assembler.\par\medskip
\noindent\textit{Editorial note.} Same-chapter pointers are located at their \textbackslash{}label lines inside the pinned chapter file; cross-chapter pointers are located in the project's own tag list staged in the snapshot (sources/stacks-chow/tags.txt), at the verbatim lines carrying their canonical tags.\par\medskip
\noindent\textit{Editorial note.} The author's item (5) of the enumeration at source line 9443, 'add more here,', is retained verbatim as part of the statement; it is the author's own placeholder in the published text and was not corrected or removed.\par\medskip
\noindent\textit{Allowed standard prerequisites (editorial).}\begin{itemize}
\item Multiplicativity of the total Chern class of a short exact sequence of finite locally free modules. Same chapter chow.tex, label lemma-additivity-chern-classes, invoked at proof line 9472. Gives c(E\_2) = c(E\_1) c(C(alpha)) once all three complexes are bounded of fixed rank, which is the last step of the argument.
\item Additivity of the Chern character. Same chapter chow.tex, label lemma-chern-character-additive, invoked at proof line 9474. Supplies the same statement for ch and for the Chern polynomials P\_p, which the proof says follow by the same argument.
\item Definiteness of the Chern classes after an envelope and existence of resolutions. Same chapter chow.tex, label lemma-chern-classes-defined, invoked at proof line 9480; Derived Categories of Schemes, label perfect-proposition-perfect-resolution-property, invoked at proof line 9487. Produces, when the irreducible components of X are quasi-compact, an envelope Y -> X splitting into quasi-compact pieces with the resolution property, so that condition (3) implies condition (1).
\end{itemize}\medskip
\ifdefined\hypertarget\hypertarget{sourcepointer-1}{}\fi
\noindent\textit{Source reference.} Original lemma 02UI: sources/stacks-chow/pinned-original-source.tex, source lines 7965--7965. Navigation pointer for \textbackslash{}ref\{lemma-additivity-chern-classes\} used inside the counted span (line 9472); the argument that consumes it is retained in the span.\par\medskip
\ifdefined\hypertarget\hypertarget{sourcepointer-2}{}\fi
\noindent\textit{Source reference.} Original lemma 0F9C: sources/stacks-chow/pinned-original-source.tex, source lines 8958--8958. Navigation pointer for \textbackslash{}ref\{lemma-chern-character-additive\} used inside the counted span (line 9474); the argument that consumes it is retained in the span.\par\medskip
\ifdefined\hypertarget\hypertarget{sourcepointer-3}{}\fi
\noindent\textit{Source reference.} Original lemma 0GUE: sources/stacks-chow/pinned-original-source.tex, source lines 9300--9300. Navigation pointer for \textbackslash{}ref\{lemma-chern-classes-defined\} used inside the counted span (line 9480); the argument that consumes it is retained in the span.\par\medskip
\ifdefined\hypertarget\hypertarget{sourcepointer-4}{}\fi
\noindent\textit{Source reference.} Original lemma 0F8I: sources/stacks-chow/tags.txt, source lines 19043--19043. Navigation pointer for \textbackslash{}ref\{perfect-proposition-perfect-resolution-property\} used inside the counted span (line 9487); the argument that consumes it is retained in the span.\par\medskip
\ifdefined\hypertarget\hypertarget{sourcepointer-5}{}\fi
\noindent\textit{Source reference.} Original lemma 02QL: sources/stacks-chow/pinned-original-source.tex, source lines 1372--1372. Navigation pointer for \textbackslash{}ref\{situation-setup\} used inside the counted span (line 9428); the argument that consumes it is retained in the span.\par\medskip
\begin{lemma}
\label{lemma-additivity-on-perfect}
Let $(S, \delta)$ be as in Situation \ref{situation-setup}.
Let $X$ be locally of finite type over $S$. Let
$$
E_1 \to E_2 \to E_3 \to E_1[1]
$$
be a distinguished triangle of perfect objects in $D(\mathcal{O}_X)$.
If one of the following conditions holds
\begin{enumerate}
\item there exists an envelope $f : Y \to X$ such that
$Lf^*E_1 \to Lf^*E_2$ can be represented by a map of locally
bounded complexes of finite locally free $\mathcal{O}_Y$-modules,
\item $E_1 \to E_2$ can be represented be a map of locally bounded complexes
of finite locally free $\mathcal{O}_X$-modules,
\item the irreducible components of $X$ are quasi-compact,
\item $X$ is quasi-compact, or
\item add more here,
\end{enumerate}
then the Chern classes of $E_1$, $E_2$, $E_3$ are defined and we have
$c(E_2) = c(E_1) c(E_3)$, $ch(E_2) = ch(E_1) + ch(E_3)$, and
$P_p(E_2) = P_p(E_1) + P_p(E_3)$.
\end{lemma}

\begin{proof}
Let $f : Y \to X$ be an envelope and let
$\alpha^\bullet : \mathcal{E}_1^\bullet \to \mathcal{E}_2^\bullet$
be a map of locally bounded complexes of finite locally free
$\mathcal{O}_Y$-modules representing $Lf^*E_1 \to Lf^*E_2$.
Then the cone $C(\alpha)^\bullet$ represents $Lf^*E_3$. Since
$C(\alpha)^n = \mathcal{E}_2^n \oplus \mathcal{E}_1^{n + 1}$
we see that $C(\alpha)^\bullet$ is a locally bounded complex
of finite locally free $\mathcal{O}_Y$-modules.
We conclude that the Chern classes of $E_1$, $E_2$, $E_3$
are defined. Moreover, recall that $c_p(E_1)$ is defined
as the unique element of $A^p(X)$ which restricts to
$c_p(\mathcal{E}_1^\bullet)$ in $A^p(Y)$. Similarly for
$E_2$ and $E_3$. Hence it suffices
to prove $c(\mathcal{E}_2^\bullet) =
c(\mathcal{E}_1^\bullet) c(C(\alpha)^\bullet)$
in $\prod_{p \geq 0} A^p(Y)$.
In turn, it suffices to prove this after restricting
to a connected component of $Y$. Hence we may assume
the complexes $\mathcal{E}_1^\bullet$ and $\mathcal{E}_2^\bullet$
are bounded complexes of finite locally free $\mathcal{O}_Y$-modules
of fixed rank. In this case the desired equality follows from the
multiplicativity of Lemma \ref{lemma-additivity-chern-classes}.
In the case of $ch$ or $P_p$ we use
Lemmas \ref{lemma-chern-character-additive}.

\medskip\noindent
In the previous paragraph we have seen that the lemma holds if
condition (1) is satisfied. Since (2) implies (1) this deals with
the second case. Assume (3). Arguing exactly as in the proof of
Lemma \ref{lemma-chern-classes-defined} we find an envelope
$f : Y \to X$ such that $Y$ is a disjoint union $Y = \coprod Y_i$
of quasi-compact (and quasi-separated) schemes each having the
resolution property. Then we may represent the restriction of
$Lf^*E_1 \to Lf^*E_2$ to $Y_i$ by a map of bounded complexes
of finite locally free modules, see
Derived Categories of Schemes, Proposition
\ref{perfect-proposition-perfect-resolution-property}.
In this way we see that condition (3) implies condition (1).
Of course condition (4) implies condition (3) and the proof
is complete.
\end{proof}

\end{document}
